\par\Needspace{6\baselineskip}\medskip
\begin{center}\hypertarget{native-179-article133}{}\pdfbookmark[2]{133.}{article-133}133.\end{center}
Quum cosinus angulorum $f,f',f''$, perinde ut quantitates $\eta,\eta''$ ab unitate
differentia secundi ordinis discrepent, patet, si pro $\dfrac{n+n''}{n'}$ valor approximatus
$1+\dfrac{\theta\theta''}{2rr'r''}$ introducatur, errorem quarti ordinis committi. Quodsi itaque loco
aequationis art. 114 haecce adhibetur
\[
a\delta'=b+\frac{c\theta+d\theta''}{\theta'}\left(1+\frac{\theta\theta''}{2rr'r''}\right)\mathpunct{\raisebox{.45ex}{,}}
\]
\hypertarget{native-180}{}in valorem distantiae $\delta'$ redundabit error secundi ordinis, quando observationes
extremae a media aequidistant, vel primi ordinis in casibus reliquis. Sed
haecce nova aequationis illius forma ad determinationem ipsius $\delta'$ haud idonea
est, quia quantitates adhuc incognitas $r,r',r''$ involvit.
\par
Iam generaliter loquendo quantitates $\dfrac{r}{r'}\mathpunct{\raisebox{.45ex}{,}}\;\dfrac{r''}{r'}$ ab unitate differentia primi
ordinis distant, et perinde etiam productum $\dfrac{rr''}{r'r'}$: in casu speciali saepius commemorato
facile perspicitur, hoc productum differentia secundi ordinis tantum
ab unitate discrepare. Quin adeo quoties orbita ellipsis parum excentrica est,
ita ut excentricitatem tamquam quantitatem primi ordinis spectare liceat, differentia
$\dfrac{rr''}{r'r'}$ ad ordinem uno gradu adhuc altiorem referri poterit. Manifestum
est itaque, errorem illum eiusdem ordinis ut antea manere, si in aequatione
nostra pro $\dfrac{\theta\theta''}{2rr'r''}$ substituatur $\dfrac{\theta\theta''}{2r'^3}\mathpunct{\raisebox{.45ex}{,}}$ unde nanciscitur formam sequentem
\[
a\delta'=b+\frac{c\theta+d\theta''}{\theta'}\left(1+\frac{\theta\theta''}{2r'^3}\right)\mathpunct{\cdot}
\]
Continet quidem haec aequatio etiamnum quantitatem incognitam $r'$, quam
tamen eliminari posse patet, quum tantummodo a $\delta'$ atque quantitatibus cognitis
pendeat. Quodsi dein aequatio rite ordinaretur, ad \textit{octavum} gradum ascenderet.
